% Checkpoint37. Complete actual based deletion proof; no checkpoint36 edit.
\statement{A.126}{Proved}{The actual three elliptic deletions and their common transporter word}
Retain the actual completed elliptic groups \(G,P\), the words
\(X,Y,\sigma,\tau\), and the unique transporter \(p_w\in P'\) of
A.121--A.125. Put
\[
 \begin{gathered}
 K=[Y,X]=\sigma\tau^2\sigma,\quad
 C=\sigma\tau=L^{-1},\quad \rho=\operatorname{Ad}(C),\\
 a=\sigma^2,\quad b=\tau^2,\quad
 m=\tau^{-1}\sigma^2\tau=K a^{-1}.
 \end{gathered}                                               \tag{A.126.1}
\]
Let \(F=\widehat F(X,Y)\) be the actual two-strand pure target.
There are specified based geometric deletion maps
\(\pi_i:P\to F\), with a specified half-twist action \(j\), satisfying
\[
 \begin{gathered}
 j(X)=X^{-1}K^{-1},\quad j(Y)=KY^{-1},\quad
 j(K)=K,\quad j^2=\operatorname{Ad}(K),\\
 \pi_1\rho=\pi_3,\qquad
 \pi_2\rho=j\pi_1,\qquad \pi_3\rho=j\pi_2.
 \end{gathered}                                               \tag{A.126.2}
\]
For every profinite commutator word \(w\) satisfying the original
pentagon, put \(U=w(X,K)\) and
\(\theta=\operatorname{Ad}(Y^{-1})j\). There is a unique
continuously determined \(a_w\in F'\) such that
\[
 \pi_1(p_w)=\pi_3(p_w)=Ua_w,\qquad
 \pi_2(p_w)=U\theta(a_w).                                    \tag{A.126.3}
\]
Moreover the entire joint derived image is
\[
             (\pi_1,\pi_2,\pi_3)(P')=(F')^3.                 \tag{A.126.4}
\]
The specified word \(a_w\) remains tied to the full transporter
equation; neither \(a_w=1\) nor realizability of every \(a\in F'\)
by an actual pentagon word is asserted.

\begin{proof}
\emph{1. The two-strand marking and based maps.}
Use the original reduced elliptic presentation
\citep[\S1.2]{A_Enriquez2014}, as fixed in A.121. In the
two-strand target its half-twist \(s\) satisfies
\[
 sXs=X^{-1},\qquad s^{-1}Ys^{-1}=Y^{-1},\qquad s^2=K.
\]
Thus \(j=\operatorname{Ad}(s)|_F\) has exactly the first row of
(A.126.2). In particular it is not replaced by an involution.
Let \(\pi_3\) be geometric deletion of strand three, with the
inherited ordering and handle paths of the two remaining strands.
Its complete row is
\[
\begin{array}{c|rrrrrr}
 &X&Y&X_2&Y_2&a&b\\ \hline
 \pi_3&X&Y&X^{-1}&Y^{-1}&K&1 .
\end{array}                                                  \tag{A.126.5}
\]
Deleting a removed handle or its meridian contracts that path;
the surviving handle products are one in the reduced two-strand
group. These facts give every entry. The ordinary full twist
\(C^3\) deletes to \(s^2=K\), so
\[
          \pi_3\rho^3=\operatorname{Ad}(K)\pi_3=j^2\pi_3.
                                                                  \tag{A.126.6}
\]
Define the other based maps, including their whiskers, by
\[
             \pi_1=\pi_3\rho^{-1},\qquad
             \pi_2=j^{-1}\pi_3\rho.                           \tag{A.126.7}
\]
These are geometric deletion of strands one and two with the
stated target basepaths, since \(C\) permutes the three labels.

\emph{2. Every original generator row.}
The pure group is generated by \(X,Y,X_2,Y_2,a,b\).
Indeed the forgetting fibration has a free rank-two base and a
punctured-torus free rank-three fiber. Its third-strand handle
generators are \(X_3=(XX_2)^{-1}\), \(Y_3=(YY_2)^{-1}\), and
a fiber meridian is \(m=Ka^{-1}\). Thus the displayed set contains
both base and fiber generators, also after full completion.
The exact conjugation table is
\[
\begin{array}{c|l}
 g&\rho(g)\\ \hline
 X&X_2a^{-1}\\
 Y&aY_2\\
 X_2&X_3m^{-1}\\
 Y_2&mY_3\\
 a&b\\
 b&m .
\end{array}                                                  \tag{A.126.8}
\]
For completeness, \(\tau\) commutes with \(X,Y\), while
\(X_2=\sigma X\sigma\) and \(Y_2=\sigma^{-1}Y\sigma^{-1}\).
These give the first two rows. The braid relation gives
\(\rho(\sigma)=\tau\), so
\[
 \rho(X_2)=\tau(X_2a^{-1})\tau=X_3m^{-1},\qquad
 \rho(Y_2)=\tau^{-1}(aY_2)\tau^{-1}=mY_3.
\]
The same braid relation gives
\[
 \rho(a)=b,\qquad
 \rho(b)=\sigma\tau^2\sigma^{-1}
        =\tau^{-1}\sigma^2\tau=m,\qquad \rho(m)=a.
                                                                  \tag{A.126.9}
\]
The last equality also follows because \(C^3\) is central in the
ordinary three-strand braid subgroup. This literal last row
corrects the earlier research table, which printed \(bmb^{-1}\);
all prior bytes are preserved and the correction is recorded
separately. It does not change any deletion row below.

Applying (A.126.5)--(A.126.9) gives the entire based table
\[
\begin{array}{c|rrrrrr}
 &X&Y&X_2&Y_2&a&b\\ \hline
 \pi_1&1&1&X&Y&1&K\\
 \pi_2&X&Y&1&1&1&1\\
 \pi_3&X&Y&X^{-1}&Y^{-1}&K&1 .
\end{array}                                                  \tag{A.126.10}
\]
Here \((\pi_1(m),\pi_2(m),\pi_3(m))=(1,K,1)\).
To check every whisker, \(\pi_3\rho\) sends \(X,Y\) to
\(j(X),j(Y)\), and sends each other generator to one, giving
the second row. For the first, use
\(\pi_3\rho^{-1}=j^{-2}\pi_3\rho^2\).
The \(\rho^2\)-images of \(X,Y,a\) have trivial \(\pi_3\)-image;
those of \(X_2,Y_2,b\) have images
\(KXK^{-1},KYK^{-1},K\), respectively. Since
\(j^{-2}=\operatorname{Ad}(K^{-1})\), the first row follows.
The identities in (A.126.2) now follow on every generator,
or directly from (A.126.6)--(A.126.7).

\emph{3. The original transporter and its full projected equations.}
Write
\[
 W=w(YX^{-1}Y^{-1},K)Yw(X,K)^{-1},\qquad B_w=w(a,b).
\]
The exact equality \(\Pi_w=p_w\Pi_0p_w^{-1}\) of A.124 is
\[
                   WB_w=\rho(p_w)Yp_w^{-1}.                 \tag{A.126.11}
\]
The three deletion images on its left are \((1,W,W)\):
\(\pi_1\) kills \(X,Y\), while each evaluation of \(B_w\)
has one argument equal to one. If \(p_i=\pi_i(p_w)\),
the full based equations are therefore
\[
       p_1=p_3=t,\qquad
       W=j(t)Yp_2^{-1}=j(p_2)Yt^{-1}.                       \tag{A.126.12}
\]
These equations retain the actual half-twist.

\emph{4. The one remaining word.}
Since \(j(X)=YX^{-1}Y^{-1}\) and \(j(K)=K\),
\[
 W=j(U)YU^{-1},\qquad
 \theta(X)=X^{-1},\quad\theta(Y)=XY^{-1}X^{-1},\quad
 \theta^2=1,\quad j\theta=\operatorname{Ad}(Y).               \tag{A.126.13}
\]
For example the inner coefficient for \(\theta^2\) is
\(Y^{-1}j(Y^{-1})K=1\). Set \(a_w=U^{-1}t\). The first
equation for \(W\) in (A.126.12) gives
\(p_2=U\theta(a_w)\); the other is then automatic by
\(j\theta=\operatorname{Ad}(Y)\). This proves (A.126.3).
A.125 gives \(p_w\in P'\), so all \(p_i\) and \(a_w\) are
in \(F'\). Uniqueness and continuity follow from the actual
unique continuous \(p_w\). The frequently derived equality
\(j(W)=KW^{-1}\) already follows from (A.126.13) and
\(j^2=\operatorname{Ad}(K)\); it adds no restriction.

\emph{5. The whole joint derived image.}
A.125's integral intersection calculation puts all three pure
pair meridians in \(P'\). Their joint images are
\[
       b\mapsto(K,1,1),\qquad m\mapsto(1,K,1),
       \qquad a\mapsto(1,1,K).
\]
Each \(\pi_i\) is onto. Therefore the image of the closed normal
closure of each meridian contains the entire normal closure of
\(K\) in its one target coordinate, and is trivial in the other
two. The normal closure of \(K=[Y,X]\) in \(F\) is exactly
\(F'\), by the universal free profinite abelian quotient.
The images are closed by compactness. Thus the joint image of
\(P'\) contains all three coordinate copies of \(F'\), and
functoriality supplies the reverse containment. This proves
(A.126.4). An arbitrary triple of the form (A.126.3) consequently
has some derived pure lift; that lift need not satisfy the
unprojected equation (A.126.11).
\end{proof}

\paragraph{Marking and scope.}
The maps above use geometric deletion and every checked generator
row, rather than an inconsistent extracted forgetting table
discussed in the retained research audit. No claim about an
uninspected printed line is required. The original boundary,
the canonical common word, and the original CM class remain
unevaluated. The full research and literal-row correction are
retained in
\path{work_A/actual_elliptic_based_forgetting36.md} and
\path{work_A/actual_based_forgetting_row_correction37.md}.
